\documentclass[../../main.tex]{subfiles}
\begin{document}

\begin{flushright}
\footnotesize\textit{【自動文字起こし・要確認】}
\end{flushright}

{\bf [解]}
\textbf{(1)} $\sqrt[3]{2}$が有理数だとする．この時$a\in\mathbb N$，$b\in\mathbb N$，$a\perp b$として
\begin{align}
\sqrt[3]{2}=\dfrac{b}{a}\quad\therefore\ 2a^3=b^3
\end{align}
とかける．したがって$b^3$が$2$の因数を持ち，$2$は素数だから$b$は$2$の倍数．この時$b=2b'\ (b'\in\mathbb N)$とかけるから
\begin{align}
a^3=4b'^3
\end{align}
となり同じく$a$も$2$の倍数となり，$a\perp b$に反する．よって矛盾し，$\sqrt[3]{2}\notin\mathbb Q$．$\blacksquare$

\textbf{(2)} $\alpha=\sqrt[3]{2}$とおく．$Q(x)$を有理数係数の整式として
\begin{align}
P(x)=Q(x)(x^3-2)+ax^2+bx+c\quad\cdots\text{※}
\end{align}
とおける．ただし，$a,b,c\in\mathbb{Q}$である．$P(\alpha)=0$から
\begin{align}
a\alpha^2+b\alpha+c=0\quad\cdots①
\end{align}

ここで，$a=0$の時，$b=c=0$ならOKだが，その他の場合は（1)から）$\alpha$が$b,c\in\mathbb Q$を用いて有理数と表せることになり矛盾．$\quad\cdots②$

又，$a\ne0$の時，$x^3-2$を$R(x)=ax^2+bx+c$で割った商$S(x)$，あまり$T(x)$として
\begin{align}
x^3-2=R(x)S(x)+T(x)\quad\cdots\text{※}
\end{align}
ただし$T(x)$は1次以下の整式で，$x=\alpha$として
\begin{align}
T(\alpha)=0
\end{align}
だが，これは上記の$a=0$の場合と同様にして，$T(x)\equiv0$となる．よって
\begin{align}
x^3-2=R(x)S(x)=(x^2+\alpha x+\alpha^2)(x-\alpha)
\end{align}
で$R(x)$が2次式だから
\begin{align}
R(x)=x^2+\alpha x+\alpha^2
\end{align}
となり，$R(x)$の係数が有理数であることに反し矛盾．$\quad\cdots③$

②③から，$a=b=c=0$で，この時$P(x)$は$x^3-2$で割り切れる．$\blacksquare$

\bigskip
\noindent{\bf [解2]}
（前略）
\begin{align}
a\alpha^2+b\alpha+c=0\quad\cdots④
\end{align}
両辺に$\alpha$をかけて，$\alpha^3=2$から
\begin{align}
b\alpha^2+c\alpha+2a=0\quad\cdots⑤
\end{align}

④$\times b$－⑤$\times a$より
\begin{align}
(b^2-ac)\alpha+bc-2a^2=0
\end{align}

$\alpha\notin\mathbb Q$だから
\begin{align}
b^2-ac=0，\qquad bc-2a^2=0
\end{align}

ここで，$c$を消去して，$a\ne0$とすると
\begin{align}
b^3=2a^3\quad\therefore\ \alpha=\dfrac{b}{a}
\end{align}
から，$\alpha\in\mathbb Q$に反し矛盾．よって$a=b=c=0$となり，示された．$\blacksquare$

\newpage

\end{document}
